%% SYNTHETIC TEST UNIT generated by writing/tools/make_outline_book.py from FINAL_OUTLINE_v2 -- not book text
\chapter{严格核心 III：部分子分布与碎裂函数的算符定义}\label{ch:pdf}
本单元为测试用的合成文本（由 FINAL\_OUTLINE\_v2 自动生成），陈述内容从略。

\section{第 1 节（测试）}\label{sec:pdf-definitions}

\begin{definition}[夸克分布]\label{def:pdf-quark}
（测试用合成陈述：内容从略。）
\uses{def:prelim-wilson, def:prelim-lightfront, def:prelim-msbar}
\end{definition}

\begin{definition}[反夸克分布]\label{def:pdf-antiquark}
（测试用合成陈述：内容从略。）
\uses{def:pdf-quark}
\end{definition}

\begin{definition}[胶子分布]\label{def:pdf-gluon}
（测试用合成陈述：内容从略。）
\uses{def:prelim-wilson, def:pdf-quark}
\end{definition}

\begin{proposition}[规范不变性、推动不变性与自由场值 \tagR]\label{prop:pdf-invariance}
（$R_{\mathrm{def}}$：由算符定义直接推出。）
（测试用合成陈述：内容从略。）
非微扰地位：\cref{asm:intro-wightman}、\cref{asm:intro-composite}。
\uses{def:pdf-quark, def:pdf-antiquark, def:pdf-gluon, prop:prelim-wilson-covariance, lem:prelim-boost, asm:intro-wightman, asm:intro-composite}
\end{proposition}

\section{第 2 节（测试）}\label{sec:pdf-properties}

\begin{example}[自由场值]\label{ex:pdf-free-field}
（测试用合成陈述：内容从略。）
\uses{prop:pdf-invariance}
\end{example}

\begin{theorem}[支撑（微扰论） \tagR]\label{thm:pdf-support-pt}
（$R_{\mathrm{PT}}$：在重整化微扰论中对红外正规化的部分子态严格。）
（测试用合成陈述：内容从略。）
（本命题只涉及部分子态。）
\uses{def:pdf-quark, def:pdf-gluon, def:ir-landau-setup}
\end{theorem}

\begin{proposition}[支撑（强子态） \tagA]\label{prop:pdf-support-hadron}
（测试用合成陈述：内容从略。）
非微扰地位：\cref{asm:intro-wightman}、\cref{asm:intro-composite}。
\uses{def:pdf-quark, def:pdf-antiquark, thm:pdf-support-pt, asm:intro-wightman, asm:intro-composite}
\end{proposition}
\begin{proofsketch}测试。\end{proofsketch}
\paragraph{地位与缺口}测试：本陈述的地位由\cref{rem:pdf-support-jaffe}决定。

\begin{remark}[支撑的第二种论证：$k^-$ 极点]\label{rem:pdf-support-jaffe}
（测试用合成陈述：内容从略。）
\uses{thm:pdf-support-pt}
\end{remark}

\section{第 3 节（测试）}\label{sec:pdf-interpretation}

\begin{remark}[数密度诠释（启发式）]\label{rem:pdf-number-density}
（测试用合成陈述：内容从略。）
% uses: none
\end{remark}

\begin{remark}[光锥规范的剩余规范自由与 $1/k^+$ 处方]\label{rem:pdf-gauge-residual}
（测试用合成陈述：内容从略。）
% uses: none
\end{remark}

\begin{definition}[光线算符]\label{def:pdf-lightray}
（测试用合成陈述：内容从略。）
\uses{def:pdf-quark}
\end{definition}

\begin{lemma}[Wilson 线顶点的矩 = 定域算符顶点 \tagR]\label{lem:pdf-wilson-moments}
（$R_{\mathrm{def}}$：由算符定义直接推出。）
（测试用合成陈述：内容从略。）
\uses{def:pdf-lightray, prop:prelim-wilson-rules, lem:appL-divided-differences}
\end{lemma}

\section{第 4 节（测试）}\label{sec:pdf-moments}

\begin{theorem}[PDF 的矩 = 扭度 2 定域算符矩阵元 \tagR]\label{thm:pdf-moments-local}
（$R_{\mathrm{PT}}$：在重整化微扰论中对红外正规化的部分子态严格。）
（测试用合成陈述：内容从略。）
非微扰地位：\cref{asm:intro-wightman}、\cref{asm:intro-composite}。
\uses{lem:pdf-wilson-moments, thm:prelim-composite-renorm, def:dis-twist, prop:appL-bare, prop:appL-renormalized, asm:intro-wightman, asm:intro-composite}
\end{theorem}

\begin{theorem}[由算符重整化得到 DGLAP 方程 \tagR]\label{thm:pdf-dglap}
（$R_{\mathrm{PT}}$：在重整化微扰论中对红外正规化的部分子态严格。）
（测试用合成陈述：内容从略。）
非微扰地位：\cref{asm:intro-composite}。
\uses{def:pdf-quark, def:pdf-gluon, thm:pdf-moments-local, prop:qcd-rge, prop:prelim-beta, asm:intro-composite}
\end{theorem}

\begin{proposition}[领头阶分裂函数 \tagR]\label{prop:pdf-lo-kernels}
（$R_{\mathrm{PT}}$：在重整化微扰论中对红外正规化的部分子态严格。）
（测试用合成陈述：内容从略。）
\uses{thm:pdf-dglap, def:pdf-quark, def:pdf-gluon}
\end{proposition}

\begin{proposition}[Feynman 规范下由算符定义计算 $P_{qq}$ \tagR]\label{prop:pdf-pqq-operator}
（$R_{\mathrm{PT}}$：在重整化微扰论中对红外正规化的部分子态严格。）
（测试用合成陈述：内容从略。）
\uses{def:pdf-quark, def:prelim-feynman-rules, prop:prelim-wilson-rules, def:prelim-msbar}
\end{proposition}

\section{第 5 节（测试）}\label{sec:pdf-dglap}

\begin{example}[单圈夸克 PDF（逐图）]\label{ex:pdf-oneloop}
（测试用合成陈述：内容从略。）
\uses{prop:pdf-pqq-operator}
\end{example}

\begin{proposition}[纯维数正规化中的部分子分布 \tagR]\label{prop:pdf-partonic-dr}
（$R_{\mathrm{PT}}$：在重整化微扰论中对红外正规化的部分子态严格。）
（测试用合成陈述：内容从略。）
\uses{def:pdf-quark, lem:prelim-scaleless, ex:prelim-scaleless, prop:pdf-lo-kernels}
\end{proposition}

\begin{proposition}[定域算符的单圈反常维度（含混合） \tagR]\label{prop:pdf-local-anomalous}
（$R_{\mathrm{PT}}$：在重整化微扰论中对红外正规化的部分子态严格。）
（测试用合成陈述：内容从略。）
\uses{def:dis-twist, thm:prelim-composite-renorm, prop:pdf-lo-kernels}
\end{proposition}

\begin{example}[精确到 $\varepsilon$ 的单圈矩]\label{ex:pdf-moments-exact}
（测试用合成陈述：内容从略。）
\uses{thm:pdf-moments-local, prop:pdf-pqq-operator}
\end{example}

\section{第 6 节（测试）}\label{sec:pdf-oneloop}

\begin{example}[胶子中的夸克]\label{ex:pdf-quark-in-gluon}
（测试用合成陈述：内容从略。）
\uses{def:pdf-gluon, prop:pdf-lo-kernels}
\end{example}

\begin{theorem}[DIS 的 $x$ 空间因子化（扭度 2 部分） \tagR]\label{thm:pdf-dis-xspace}
（$R_{\mathrm{OPE}}$：以 OPE 为前提而严格。QCD 中含夸克的微扰 OPE 是全阶论证（\cref{thm:dis-ope}），严格证明只对标量理论与纯 Yang--Mills 的规范不变算符给出；对强子态还须非微扰 OPE（\cref{asm:intro-np-ope}）。）
（测试用合成陈述：内容从略。）
非微扰地位：\cref{asm:intro-wightman}、\cref{asm:intro-composite}、\cref{asm:intro-np-ope}、\cref{asm:intro-af-coeff}、\cref{asm:intro-growth}。
\uses{thm:dis-moments, thm:pdf-moments-local, lem:prelim-mellin, cor:appK-pdf, prop:pdf-partonic-dr, asm:intro-wightman, asm:intro-composite, asm:intro-np-ope, asm:intro-af-coeff, asm:intro-growth}
\end{theorem}

\begin{proposition}[求和规则 \tagR]\label{prop:pdf-sum-rules}
（$R_{\mathrm{PT}}$：在重整化微扰论中对红外正规化的部分子态严格。）
（测试用合成陈述：内容从略。）
非微扰地位：\cref{asm:intro-wightman}、\cref{asm:intro-composite}、\cref{asm:intro-growth}、\cref{prop:dis-low-moments}。
\uses{thm:pdf-moments-local, prop:pdf-lo-kernels, thm:pdf-dglap, asm:intro-wightman, asm:intro-composite}
\end{proposition}

\begin{remark}[解析延拓与 $P^\pm_{\rm NS}$]\label{rem:pdf-ns-plusminus}
（测试用合成陈述：内容从略。）
% uses: none
\end{remark}

\section{第 7 节（测试）}\label{sec:pdf-dis}

\begin{remark}[方案变换]\label{rem:pdf-schemes}
（测试用合成陈述：内容从略。）
% uses: none
\end{remark}

\begin{conjecture}[$\overline{\rm MS}$ PDF 的正定性 \tagC]\label{conj:pdf-positivity}
（测试用合成陈述：内容从略。）
\uses{def:pdf-quark, def:pdf-gluon}
\end{conjecture}
\paragraph{地位与缺口}测试：见文献。

\begin{definition}[碎裂函数]\label{def:pdf-ff}
（测试用合成陈述：内容从略。）
\uses{def:prelim-wilson, def:prelim-lightfront}
\end{definition}

\begin{remark}[碎裂函数的乘法可重整性是独立输入]\label{rem:pdf-ff-multren}
（测试用合成陈述：内容从略。）
\uses{def:pdf-ff, thm:pdf-moments-local}
\end{remark}

\section{第 8 节（测试）}\label{sec:pdf-sumrules}

\begin{proposition}[类时 DGLAP 演化 \tagA]\label{prop:pdf-timelike-dglap}
（测试用合成陈述：内容从略。）
非微扰地位：\cref{asm:intro-wightman}、\cref{asm:intro-composite}。
\uses{def:pdf-ff, rem:pdf-ff-multren, prop:pdf-lo-kernels, asm:appN-multren, asm:intro-wightman, asm:intro-composite}
\end{proposition}
\begin{proofsketch}测试。\end{proofsketch}
\paragraph{地位与缺口}测试：本陈述的地位由\cref{asm:appN-multren}决定。

\begin{definition}[螺旋度与横度分布]\label{def:pdf-polarized}
（测试用合成陈述：内容从略。）
\uses{def:pdf-quark, def:prelim-gamma5}
\end{definition}
